\documentclass[12pt,a4paper]{article}
\usepackage[margin=1in]{geometry}
\usepackage[T1]{fontenc}
\usepackage[utf8]{inputenc}
\usepackage{lmodern}
\usepackage{hyperref}
\usepackage{longtable}
\usepackage{array}
\usepackage{booktabs}
\usepackage{enumitem}
\usepackage{xcolor}
\hypersetup{
  colorlinks=true,
  linkcolor=blue,
  urlcolor=blue,
  pdftitle={Medication Prescription Tracker Analysis Report}
}

\title{Medication Prescription Tracker\\Analysis Report}
\author{Generated from repository inspection}
\date{March 17, 2026}

\begin{document}
\maketitle
\tableofcontents
\newpage

\section{Project Overview}
The Medication Prescription Tracker is a role-based full-stack healthcare application designed to manage prescription creation, medication schedules, patient adherence, pharmacy inventory, dispensing operations, and administrative oversight.

The implementation is split into two main subsystems:
\begin{itemize}[leftmargin=*]
  \item \textbf{Frontend:} React + Vite application in \texttt{mediui/}
  \item \textbf{Backend:} Spring Boot application in \texttt{medimanager/}
\end{itemize}

The system supports four principal roles:
\begin{itemize}[leftmargin=*]
  \item \textbf{Patient} - views prescriptions, tracks doses, reviews adherence analytics
  \item \textbf{Doctor} - creates patients, issues and manages prescriptions
  \item \textbf{Pharmacist} - manages stock and dispenses medication
  \item \textbf{Admin} - monitors users, prescriptions, and audit activity
\end{itemize}

\section{Objectives of the System}
From the implementation, the major system goals appear to be:
\begin{enumerate}[leftmargin=*]
  \item digitize prescription creation and renewal,
  \item connect prescriptions with medication schedules,
  \item measure patient adherence through dose logging,
  \item support pharmacist inventory and dispensing workflows,
  \item enforce role-based access control,
  \item provide administrative visibility over users and system activity.
\end{enumerate}

\section{Technology Stack}
\subsection{Frontend}
\begin{itemize}[leftmargin=*]
  \item React 18
  \item React Router DOM
  \item Vite
  \item Recharts for analytics visualizations
  \item Lucide React for icons
\end{itemize}

\subsection{Backend}
\begin{itemize}[leftmargin=*]
  \item Spring Boot
  \item Spring Web
  \item Spring Data JPA
  \item Spring Security
  \item JWT via \texttt{jjwt}
  \item Flyway for schema migration
  \item H2 and PostgreSQL support
  \item iText 7 for prescription PDF generation
\end{itemize}

\subsection{Programming Environment}
\begin{itemize}[leftmargin=*]
  \item Java 21
  \item JavaScript / JSX
  \item Maven for backend build
  \item npm for frontend build
\end{itemize}

\section{High-Level Architecture}
The system follows a conventional client-server design.

\subsection{Frontend Layer}
The React application handles:
\begin{itemize}[leftmargin=*]
  \item authentication state persistence,
  \item route protection based on user role,
  \item role-specific dashboards,
  \item forms for prescription, patient, inventory, and profile management,
  \item analytics visualization for patient adherence.
\end{itemize}

\subsection{Backend Layer}
The Spring Boot application is organized into:
\begin{itemize}[leftmargin=*]
  \item \textbf{Controllers} for REST endpoints,
  \item \textbf{Services} for business logic,
  \item \textbf{Repositories} for data access,
  \item \textbf{Entities} for relational persistence,
  \item \textbf{Security components} for JWT authentication and role enforcement.
\end{itemize}

\subsection{Persistence Layer}
JPA entities model users, prescriptions, prescription medicines, schedules, dose logs, inventory records, dispense records, profiles, and audit logs.

\section{Functional Modules}
\subsection{Authentication and Profile Management}
Implemented capabilities include:
\begin{itemize}[leftmargin=*]
  \item signup,
  \item login,
  \item password change,
  \item forgot password request,
  \item password reset,
  \item profile editing.
\end{itemize}

Authentication data is stored in the frontend local storage under \texttt{medimanager\_auth}. Backend security uses stateless JWT-based authorization.

\subsection{Prescription Management}
Doctors can:
\begin{itemize}[leftmargin=*]
  \item create prescriptions,
  \item renew existing prescriptions,
  \item update existing prescriptions,
  \item delete prescriptions,
  \item download prescriptions as PDF,
  \item view their own prescription list,
  \item fetch prescriptions for specific patients.
\end{itemize}

The backend automatically converts prescription medicine entries into medication schedules for the patient.

\subsection{Patient Management}
Doctors can onboard patients directly by creating:
\begin{itemize}[leftmargin=*]
  \item a user account with role \texttt{PATIENT},
  \item an associated patient profile containing demographic fields.
\end{itemize}

\subsection{Medication Scheduling and Adherence}
Patients can:
\begin{itemize}[leftmargin=*]
  \item view their prescriptions,
  \item create medication schedules,
  \item retrieve adherence percentage,
  \item access dose logs for a date range,
  \item use a medication tracker UI.
\end{itemize}

Adherence is computed from taken vs missed dose logs.

\subsection{Pharmacy Inventory and Dispensing}
Pharmacists can:
\begin{itemize}[leftmargin=*]
  \item add inventory items,
  \item update stock quantity,
  \item delete inventory items,
  \item retrieve low-stock alerts,
  \item retrieve expiring-stock alerts,
  \item view active prescriptions,
  \item dispense medication against inventory.
\end{itemize}

\subsection{Administrative Oversight}
Admins can:
\begin{itemize}[leftmargin=*]
  \item view system summary statistics,
  \item list all users,
  \item filter users by role,
  \item create users,
  \item toggle user active status,
  \item delete users,
  \item view prescriptions,
  \item inspect audit logs.
\end{itemize}

\section{Comprehensive Feature Inventory}
This section lists the implemented features in a more exhaustive form so the document can be used for complete project analysis.

\subsection{Authentication Features}
\begin{itemize}[leftmargin=*]
  \item user signup with role-aware onboarding flow,
  \item user login with JWT token generation,
  \item persistent authenticated session in frontend local storage,
  \item protected frontend routes,
  \item backend route authorization by role,
  \item change password for authenticated users,
  \item forgot password request handling,
  \item reset password endpoint,
  \item logout support,
  \item role display in navigation and workspace routing.
\end{itemize}

\subsection{Profile and Account Features}
\begin{itemize}[leftmargin=*]
  \item profile editing for authenticated users,
  \item user activation and deactivation,
  \item admin-side user creation,
  \item admin-side user deletion,
  \item user filtering by role,
  \item account creation origin tracking through \texttt{createdBy}.
\end{itemize}

\subsection{Doctor Features}
\begin{itemize}[leftmargin=*]
  \item patient onboarding by doctor,
  \item patient search and listing,
  \item prescription creation for selected patient,
  \item multi-medicine prescription entry,
  \item diagnosis capture,
  \item expiry date control,
  \item prescription renewal,
  \item prescription update,
  \item prescription deletion,
  \item viewing doctor-owned prescriptions,
  \item viewing prescriptions for a specific patient,
  \item monitoring patient adherence,
  \item viewing patient schedules,
  \item viewing patient dose logs,
  \item generating downloadable prescription PDF files.
\end{itemize}

\subsection{Prescription Features}
\begin{itemize}[leftmargin=*]
  \item prescription issue date recording,
  \item prescription expiry date recording,
  \item prescription status management,
  \item support for active, expired, renewed, and cancelled lifecycle states,
  \item medicine-level prescription breakdown,
  \item dosage capture,
  \item timing or frequency capture,
  \item duration capture,
  \item instruction or notes capture,
  \item prescription response DTO mapping for frontend use,
  \item PDF export for prescription documentation.
\end{itemize}

\subsection{Medication Scheduling Features}
\begin{itemize}[leftmargin=*]
  \item automatic medication schedule generation from prescription medicines,
  \item manual medication schedule creation,
  \item schedule start and end date tracking,
  \item medicine time-of-day tracking,
  \item schedule activation state,
  \item daily, weekly, and monthly-style frequency handling,
  \item parsing of medicine timing into schedule time,
  \item schedule retrieval by patient,
  \item expired schedule deactivation logic,
  \item schedule deletion through deactivation.
\end{itemize}

\subsection{Patient Features}
\begin{itemize}[leftmargin=*]
  \item personal prescription dashboard,
  \item prescription detail modal view,
  \item prescription PDF download,
  \item medication tracker access,
  \item personal adherence retrieval,
  \item dose log retrieval over date ranges,
  \item visual adherence dashboard,
  \item daily taken versus missed medication chart,
  \item adherence trend chart,
  \item active prescription count display,
  \item insight and recommendation panel.
\end{itemize}

\subsection{Dose Logging and Adherence Features}
\begin{itemize}[leftmargin=*]
  \item taken and missed dose states,
  \item timestamped dose logging,
  \item optional notes with dose logs,
  \item dose retrieval by schedule,
  \item dose retrieval by patient,
  \item adherence percentage calculation,
  \item analytics refresh support in frontend after dose logging,
  \item date-range-based analytics requests from frontend.
\end{itemize}

\subsection{Pharmacist Features}
\begin{itemize}[leftmargin=*]
  \item pharmacist-specific inventory dashboard,
  \item adding new drug inventory items,
  \item batch number tracking,
  \item expiry date tracking,
  \item stock quantity tracking,
  \item low-stock threshold tracking,
  \item manufacturer capture,
  \item unit price capture,
  \item update stock operation,
  \item inventory deletion,
  \item low-stock alert retrieval,
  \item expiring-soon stock retrieval,
  \item dashboard summary statistics,
  \item active prescription viewing for dispense decisions,
  \item patient-specific active prescription retrieval,
  \item medication dispensing with quantity validation,
  \item optional prescription linkage during dispensing,
  \item dispense record creation,
  \item total stock value calculation.
\end{itemize}

\subsection{Admin Features}
\begin{itemize}[leftmargin=*]
  \item system summary dashboard,
  \item user counts,
  \item doctor counts,
  \item patient counts,
  \item prescription counts,
  \item recent prescription display,
  \item full user listing,
  \item role-filtered user listing,
  \item admin-side account creation,
  \item toggle user active status,
  \item delete user,
  \item prescription oversight,
  \item audit log retrieval by user,
  \item audit log retrieval by entity,
  \item audit log retrieval by date window.
\end{itemize}

\subsection{Audit and Monitoring Features}
\begin{itemize}[leftmargin=*]
  \item audit logging for prescription creation,
  \item audit logging for prescription renewal,
  \item audit logging for prescription update,
  \item audit logging for prescription deletion,
  \item audit logging for schedule creation,
  \item audit logging for dose logging,
  \item audit visibility for administrators.
\end{itemize}

\subsection{Frontend Experience Features}
\begin{itemize}[leftmargin=*]
  \item role-based top navigation,
  \item workspace tiles on dashboard,
  \item protected route component,
  \item login page,
  \item signup page,
  \item forgot password page,
  \item profile editor component,
  \item password change component,
  \item doctor prescription management page,
  \item patient prescription page,
  \item patient analytics page,
  \item pharmacist inventory page,
  \item admin dashboard page,
  \item admin user management page.
\end{itemize}

\subsection{Backend Infrastructure Features}
\begin{itemize}[leftmargin=*]
  \item stateless security filter chain,
  \item JWT request filter,
  \item custom user details service,
  \item BCrypt password encoding,
  \item Flyway migration support,
  \item H2 development database support,
  \item PostgreSQL runtime support,
  \item REST API layering through controller-service-repository structure,
  \item PDF generation through iText,
  \item environment-based database configuration support.
\end{itemize}

\section{Role-Based Workflow Analysis}
\subsection{Doctor Workflow}
\begin{enumerate}[leftmargin=*]
  \item Create or select a patient.
  \item Enter diagnosis and prescription expiry date.
  \item Add one or more medicines with dosage, frequency, duration, and notes.
  \item Submit prescription to backend.
  \item Backend saves prescription and medicines.
  \item Backend auto-generates medication schedules.
  \item Doctor can later renew, edit, delete, or export PDF.
\end{enumerate}

\subsection{Patient Workflow}
\begin{enumerate}[leftmargin=*]
  \item Log in and open patient dashboard.
  \item Review prescription list and details.
  \item Open medication tracker and log medicine-taking behavior.
  \item View analytics dashboard for adherence trends.
  \item Download prescription PDF when needed.
\end{enumerate}

\subsection{Pharmacist Workflow}
\begin{enumerate}[leftmargin=*]
  \item Add stock items with batch, expiry, threshold, unit price, and manufacturer.
  \item Monitor low-stock and expiry alerts.
  \item Review active prescriptions.
  \item Select inventory item and dispense quantity to patient.
  \item Backend reduces stock and stores dispense record.
\end{enumerate}

\subsection{Admin Workflow}
\begin{enumerate}[leftmargin=*]
  \item Open dashboard summary.
  \item Inspect total users, patients, doctors, and prescriptions.
  \item Manage users by creation, activation toggle, or deletion.
  \item Review audit logs and prescription history.
\end{enumerate}

\subsection{Cross-Module Workflow}
\begin{enumerate}[leftmargin=*]
  \item A doctor creates a patient account and patient profile.
  \item The doctor creates a prescription with one or more medicines.
  \item The system stores the prescription and its medicine items.
  \item The backend auto-generates medication schedules from prescribed medicines.
  \item The patient views prescriptions and follows the medication tracker.
  \item Dose logs are stored as taken or missed events.
  \item Adherence analytics are calculated from patient dose history.
  \item The pharmacist reviews active prescriptions and manages stock.
  \item The pharmacist dispenses medication and inventory is updated.
  \item The admin monitors the platform through user and audit oversight.
\end{enumerate}

\section{Data Model Summary}
\subsection{Core Entities}
\begin{longtable}{>{\raggedright\arraybackslash}p{0.22\textwidth} >{\raggedright\arraybackslash}p{0.70\textwidth}}
\toprule
\textbf{Entity} & \textbf{Purpose and Main Fields} \\
\midrule
\endhead
User & Stores credentials and role assignment. Key fields: id, email, password, role, active, createdAt, updatedAt, createdBy. \\
Prescription & Connects patient and doctor with diagnosis, issued date, expiry date, status, and medicine list. \\
PrescriptionMedicine & Stores medicine lines within a prescription including medicine name, dosage, frequency, duration, and instructions. \\
MedicationSchedule & Stores reminder-oriented schedule data for a patient including medicine name, time of day, frequency, date range, and active flag. \\
DoseLog & Stores whether a dose was taken or missed, plus timestamp and notes. \\
DrugInventory & Stores pharmacist-owned stock records including drug name, batch number, expiry date, stock quantity, threshold, unit price, and manufacturer. \\
MedicationDispense & Stores dispensing events linked to inventory, pharmacist, optional prescription, patient name, quantity, and dispensed time. \\
PatientProfile & Stores patient demographic details such as name, age, and gender. \\
DoctorProfile & Stores doctor metadata such as name, specialization, and license number. \\
PharmacistProfile & Stores pharmacist profile information. \\
AuditLog & Stores tracked system actions for accountability. \\
\bottomrule
\end{longtable}

\subsection{Entity Relationships}
\begin{itemize}[leftmargin=*]
  \item One patient can have many prescriptions.
  \item One doctor can issue many prescriptions.
  \item One prescription can contain many prescription medicines.
  \item One patient can have many medication schedules.
  \item One medication schedule can have many dose logs.
  \item One pharmacist can own many inventory records.
  \item One inventory record can have many dispense events.
  \item One admin or service action can create many audit log entries over time.
\end{itemize}

\section{API Surface Summary}
\subsection{Authentication}
\begin{itemize}[leftmargin=*]
  \item \texttt{POST /api/auth/signup}
  \item \texttt{POST /api/auth/login}
  \item \texttt{PUT /api/auth/change-password}
  \item \texttt{POST /api/auth/forgot-password}
  \item \texttt{POST /api/auth/reset-password}
\end{itemize}

\subsection{Patient}
\begin{itemize}[leftmargin=*]
  \item \texttt{GET /api/patient/prescriptions}
  \item \texttt{GET /api/patient/prescriptions/\{id\}}
  \item \texttt{POST /api/patient/schedules/create}
  \item \texttt{GET /api/patient/adherence}
  \item \texttt{GET /api/patient/dose-logs}
\end{itemize}

\subsection{Doctor}
\begin{itemize}[leftmargin=*]
  \item \texttt{POST /api/doctor/prescriptions/create}
  \item \texttt{GET /api/doctor/prescriptions}
  \item \texttt{POST /api/doctor/prescriptions/\{id\}/renew}
  \item \texttt{PUT /api/doctor/prescriptions/\{id\}}
  \item \texttt{DELETE /api/doctor/prescriptions/\{id\}}
  \item \texttt{POST /api/doctor/patients}
  \item \texttt{GET /api/doctor/patients}
  \item \texttt{GET /api/doctor/patients/\{patientId\}/prescriptions}
  \item \texttt{GET /api/doctor/patients/\{patientId\}/adherence}
  \item \texttt{GET /api/doctor/patients/\{patientId\}/schedules}
  \item \texttt{GET /api/doctor/patients/\{patientId\}/dose-logs}
\end{itemize}

\subsection{Pharmacist}
\begin{itemize}[leftmargin=*]
  \item \texttt{POST /api/pharmacist/inventory/add}
  \item \texttt{GET /api/pharmacist/inventory}
  \item \texttt{PUT /api/pharmacist/inventory/\{drugId\}/update-stock}
  \item \texttt{DELETE /api/pharmacist/inventory/\{drugId\}}
  \item \texttt{GET /api/pharmacist/inventory/low-stock}
  \item \texttt{GET /api/pharmacist/inventory/expiring}
  \item \texttt{GET /api/pharmacist/prescriptions}
  \item \texttt{GET /api/pharmacist/prescriptions/patient/\{patientId\}}
  \item \texttt{GET /api/pharmacist/dashboard/stats}
  \item \texttt{POST /api/pharmacist/inventory/dispense}
\end{itemize}

\subsection{Admin}
\begin{itemize}[leftmargin=*]
  \item \texttt{GET /api/admin/summary}
  \item \texttt{GET /api/admin/prescriptions}
  \item \texttt{GET /api/admin/dashboard/stats}
  \item \texttt{GET /api/admin/users}
  \item \texttt{GET /api/admin/users/role/\{role\}}
  \item \texttt{POST /api/admin/users}
  \item \texttt{PUT /api/admin/users/\{userId\}/toggle-status}
  \item \texttt{DELETE /api/admin/users/\{userId\}}
  \item \texttt{GET /api/admin/audit-logs}
  \item \texttt{GET /api/admin/audit-logs/entity/\{entityType\}/\{entityId\}}
\end{itemize}

\subsection{General Prescription Endpoints}
\begin{itemize}[leftmargin=*]
  \item \texttt{POST /api/prescriptions/create}
  \item \texttt{GET /api/prescriptions}
  \item \texttt{GET /api/prescriptions/\{patientId\}}
  \item \texttt{GET /api/prescriptions/detail/\{id\}}
  \item \texttt{GET /api/prescriptions/\{id\}/download-pdf}
\end{itemize}

\section{Security Analysis}
\subsection{Strengths}
\begin{itemize}[leftmargin=*]
  \item Stateless JWT-based authentication is present.
  \item Password hashing uses BCrypt.
  \item Route-level and method-level role restrictions are applied.
  \item Role segregation is explicit for admin, doctor, patient, and pharmacist APIs.
\end{itemize}

\subsection{Observed Limitations}
\begin{itemize}[leftmargin=*]
  \item Some controllers expose broad \texttt{@CrossOrigin(origins = "*")} configuration.
  \item Password reset flow is placeholder-style and does not show tokenized email verification.
  \item Several endpoints return generic bad-request responses without structured error taxonomy.
  \item Public prescription endpoints coexist with role-restricted prescription endpoints, which may require further review for least-privilege design.
\end{itemize}

\section{Implementation Quality Analysis}
\subsection{Positive Design Choices}
\begin{itemize}[leftmargin=*]
  \item Clear separation between controllers, services, repositories, and entities.
  \item Role-specific frontend routing aligns with backend authorization structure.
  \item Automatic schedule creation from prescriptions improves workflow continuity.
  \item PDF export adds practical value for real-world prescription usage.
  \item Adherence analytics connect operational data with patient outcome tracking.
\end{itemize}

\subsection{Technical Observations}
\begin{itemize}[leftmargin=*]
  \item The system contains both global prescription endpoints and doctor-specific prescription endpoints, which suggests iterative expansion rather than strict consolidation.
  \item Some older UI artifacts and encoded characters indicate mixed documentation or encoding history.
  \item Adherence calculation uses aggregate taken/missed counts and does not appear to apply the provided date range at repository-count level.
  \item Update schedule logic exists but currently does not apply detailed field mutations in the shown service implementation.
  \item Frontend API access relies on local storage token retrieval, which is functional but should be evaluated for long-term security posture.
\end{itemize}

\subsection{Implemented Feature Depth}
The codebase goes beyond a basic prescription CRUD application. It includes:
\begin{itemize}[leftmargin=*]
  \item account lifecycle support,
  \item role-governed workspaces,
  \item prescription-to-schedule automation,
  \item patient adherence analytics,
  \item inventory risk monitoring,
  \item dispensing records,
  \item administrative supervision,
  \item auditability,
  \item downloadable documentation in PDF form.
\end{itemize}

\section{Analytical Strengths of the Project}
This project is suitable for academic or technical analysis because it includes:
\begin{itemize}[leftmargin=*]
  \item multi-role user management,
  \item end-to-end healthcare workflow coverage,
  \item traceable entity relationships,
  \item measurable adherence outcomes,
  \item stock and dispensing operations,
  \item auditability and administrative supervision.
\end{itemize}

These characteristics make it useful for analysis in software engineering, health informatics, information systems design, and applied database modeling.

\section{Potential Research or Evaluation Angles}
The project can be analyzed under several themes:
\begin{enumerate}[leftmargin=*]
  \item \textbf{Software architecture:} modularity, maintainability, separation of concerns.
  \item \textbf{Security:} JWT handling, role enforcement, password lifecycle, CORS policy.
  \item \textbf{Database design:} normalization, relationship mapping, transaction boundaries.
  \item \textbf{Healthcare workflow digitization:} prescription-to-schedule automation.
  \item \textbf{Analytics quality:} adherence metric design and visualization usefulness.
  \item \textbf{Usability:} role-based dashboards and workflow fit for end users.
  \item \textbf{Scalability:} readiness for production deployment and concurrent usage.
\end{enumerate}

\section{Limitations and Improvement Opportunities}
\begin{itemize}[leftmargin=*]
  \item Add stronger password reset with token issuance and email delivery.
  \item Refine CORS configuration for production environments.
  \item Consolidate duplicate or overlapping endpoint patterns.
  \item Improve validation and consistent error responses across APIs.
  \item Expand test coverage beyond the minimal visible automated test footprint.
  \item Improve adherence calculations to fully respect selected date windows.
  \item Add stronger domain rules for prescription expiration, cancellation, and dispense verification.
\end{itemize}

\section{Conclusion}
The Medication Prescription Tracker is a substantial full-stack healthcare management application with meaningful support for prescriptions, schedules, adherence, inventory, dispensing, and administration. Its architecture is practical and understandable, making it a strong candidate for system analysis. The codebase demonstrates clear functional breadth and reasonable modular organization, while still leaving identifiable opportunities for refinement in validation, security hardening, and analytical accuracy.

\section{Suggested Appendix Items for Further Study}
If a longer academic report is needed, the following appendix materials can be added later:
\begin{itemize}[leftmargin=*]
  \item UML use case diagram,
  \item entity relationship diagram,
  \item sequence diagram for prescription creation,
  \item sequence diagram for medicine dispensing,
  \item test-case matrix,
  \item security threat model,
  \item deployment architecture diagram.
\end{itemize}

\end{document}
